\chapter{Sistemas de frontera, rango dimensional y funcional de acción}
\label{chap:pantalla}

En Mecánica Dimensional, la dimensión se obtiene del número de canales
independientes disponibles en una frontera. No se añade como etiqueta externa
al espacio de estados: se calcula como el rango de un módulo de prolongaciones
sujeto a sus relaciones. Un cambio de dimensión corresponde entonces a una
variación de rango durante el transporte.

\section{Sistemas de restricción y datos de frontera}

Sea
\[
 \cdots\longrightarrow\mathcal S_{n+1}
 \xrightarrow{r_{n+1,n}}\mathcal S_n\longrightarrow\cdots
\]
un sistema inverso de estados supervivientes. La estructura de frontera en el
nivel \(n\) es la pareja
\[
 \mathcal F_n=(r_{n+1,n},\beta_n),
\]
donde \(\beta_n\) conserva el dato mínimo necesario para decidir la
compatibilidad de una prolongación. Sus coordenadas pueden incluir el tipo
aritmético, la orientación, el cociclo de acarreo y la clase de incidencia.

\begin{lemma}[Compatibilidad de las restricciones sucesivas]
Supongamos
\[
 r_{n+2,n}=r_{n+1,n}\circ r_{n+2,n+1},
 \qquad
 \beta_n\circ r_{n+1,n}
 =b_{n+1,n}\circ\beta_{n+1}.
\]
Entonces toda condición de compatibilidad expresada mediante \(\beta_{n+1}\)
induce una condición compatible en el nivel \(n\).
\end{lemma}
\begin{proof}
Para \(x\in\mathcal S_{n+2}\), la igualdad de composiciones identifica la
restricción directa con las dos restricciones consecutivas. La segunda
igualdad transporta el dato de frontera a través del mismo diagrama. La
condición en el nivel \(n\) es, por tanto, la imagen de la condición más
profunda.
\end{proof}

\section{Módulo de canales y rango dimensional}

Fijado un cuerpo \(\Kfield_n\), sea \(M_n(x)\) el módulo generado por los
canales de prolongación de \(x\), y sea \(\mathcal R_n(x)\) el submódulo de
relaciones entre ellos.

\begin{definition}[Pantalla dimensional]
La pantalla asociada a \(x\) en el nivel \(n\) es
\[
 \Sigma_n(x)=
 (\Kfield_n,M_n(x),\beta_n(x),\mathcal R_n(x)),
\]
y su dimensión efectiva es
\[
 d_{\Sigma_n}(x)=
 \dim_{\Kfield_n}\frac{M_n(x)}{\mathcal R_n(x)}.
\]
Si \(M_n(x)\cong\Kfield_n^{g_n}\) y las columnas de una matriz
\(R_n(x)\) generan \(\mathcal R_n(x)\), entonces
\[
 d_{\Sigma_n}(x)=g_n-\rank_{\Kfield_n}R_n(x).
\]
Cuando se fija un complemento y \(P_n(x)\) denota el proyector cuya imagen es
isomorfa al cociente, la misma dimensión es
\(d_{\Sigma_n}(x)=\rank_{\Kfield_n}P_n(x)\).
\end{definition}

\begin{proposition}[Invariancia por equivalencia de presentaciones]
La dimensión de pantalla es invariante bajo cambios de base invertibles en
los generadores y en las relaciones.
\end{proposition}
\begin{proof}
Un cambio de bases transforma la matriz de relaciones \(R\) en \(URV\), con
\(U\) y \(V\) invertibles. Como \(\rank(URV)=\rank R\), el número
\(g-\rank R\) no cambia. Equivalentemente, dos proyectores que presentan el
mismo cociente tienen imágenes isomorfas y, por tanto, el mismo rango.
\end{proof}

Para un transporte \(T_\gamma:\mathcal S_n\to\mathcal S_m\), definimos
\[
 \Delta d(\gamma;x)=
 d_{\Sigma_m}(T_\gamma x)-d_{\Sigma_n}(x).
\]
Una transición dimensional transporta simultáneamente el estado y su módulo
de canales; la dinámica de rango fijo es el caso \(\Delta d=0\).

\section{Pantallas aritméticas y reticulares heredadas de HMT}

Los rangos que intervienen más adelante no son etiquetas añadidas a las
ecuaciones. Proceden de construcciones lineales distintas y conservan sus
mapas de selección. El código ternario de Golay perforado
\[
 G_{11}=[11,6,5]_3\subset\mathbb F_3^{11}
\]
es perfecto de radio dos \cite{MacWilliamsSloane1977}, pues las bolas
centradas en sus \(3^6\) palabras son
disjuntas y
\[
 V_3(11,2)=1+2\binom{11}{1}+2^2\binom{11}{2}
 =1+22+220=243,
\qquad
 3^6\cdot243=3^{11}.
\]
Aquí \(243\) es el cardinal del espacio de síndromes. No se identifica por
mera igualdad numérica con la capa de \(243\) celdas de la completación
cúbica de APP\@.

La segunda cadena parte del módulo de permutación
\(V_{12}=\mathbb R^{12}\) de \(A_5\) sobre los doce vértices del icosaedro.
Para fijar el calibre, hacemos actuar \(A_5\) sobre
\(\{0,1,2,3,4\}\), tomamos
\(H=\langle(0\ 1\ 2\ 3\ 4)\rangle\cong C_5\) e identificamos la posición
\(j\) con el coset izquierdo \(r_jH\). En notación de imágenes, los doce
representantes son
\[
\begin{array}{c|c@{\qquad}c|c@{\qquad}c|c}
j&r_j&j&r_j&j&r_j\\ \hline
1&(3,2,4,1,0)&5&(3,2,0,4,1)&9&(3,1,2,4,0)\\
2&(3,1,0,2,4)&6&(4,0,1,2,3)&10&(2,1,4,3,0)\\
3&(0,1,3,4,2)&7&(4,3,2,1,0)&11&(3,4,1,2,0)\\
4&(1,4,2,3,0)&8&(0,2,3,1,4)&12&(4,2,1,3,0).
\end{array}
\]
Denotamos por \(5A\) la clase de conjugación que contiene al generador
\(g_5=(0\ 1\ 2\ 3\ 4)\) de \(H\), y por \(5B\) la clase que contiene a
\(g_5^2\). En el orden de clases \(1A,2A,3A,5A,5B\), el carácter de permutación es
\(\chi_{12}=(12,0,0,2,2)\). Para
\(\phi=(1+\sqrt5)/2\) y \(\phi'=(1-\sqrt5)/2\), elegimos el carácter
irreducible
\[
 \chi_3=(3,-1,0,\phi,\phi').
\]
El producto escalar de caracteres da la descomposición exacta
\[
 V_{12}\cong V_1\oplus V_3\oplus V_{3'}\oplus V_5.
\]

El proyector
\[
 P_{11}=I_{12}-\frac1{12}J_{12}
\]
elimina exclusivamente el vector uniforme y tiene imagen
\[
 V_{11}=\mathbf1^\perp\cong V_3\oplus V_{3'}\oplus V_5.
\]
Para el estado dodecafásico
\[
 K=(234,543,140,729,659,824,621,58,914,794,146,601),
 \qquad \sum_jK_j=6263,
\]
sea \(\widetilde K=P_{11}K\). La correspondencia posición--coset anterior
define las matrices de permutación \(\rho(g)\), y el proyector central
\[
 P_3=\frac3{60}\sum_{g\in A_5}\chi_3(g^{-1})\rho(g)
\]
satisface \(P_3^2=P_3=P_3^{\mathsf T}\) y tiene rango tres. La evaluación
exacta en \(\mathbb Q(\sqrt5)\) produce
\[
 u_K=P_3\widetilde K,
 \qquad
 \|u_K\|^2=\frac{6638585+2275584\sqrt5}{20}>0.
\]
Por consiguiente, la recta \(\ell_K=\mathbb Ru_K\) está bien definida y
\[
 V_{10}:=(V_3\cap\ell_K^\perp)\oplus V_{3'}\oplus V_5,
 \qquad
 V_{11}=\ell_K\oplus V_{10},
 \qquad \dim V_{10}=10.
\]
La cadena \(12\to11\to10\) significa, respectivamente, eliminar la
componente uniforme y la recta polarizada por \(K\). El conjunto finito de
los sesenta elementos se divide en doce clases; su enumeración
reconstruye el proyector y la norma exactamente.

La reducción \(26\to24\) pertenece a otra categoría. Sea \(U\) el plano
hiperbólico entero con matriz de Gram
\(\left(\begin{smallmatrix}0&1\\1&0\end{smallmatrix}\right)\), y sea
\[
 L_{26}=\Lambda_{24}\oplus U.
\]
Con el sumando hiperbólico fijado, ésta es la presentación estándar del
retículo par unimodular lorentziano \(II_{25,1}\)
\cite{ConwaySloane1999}.
Para un vector isotrópico primitivo \(e_+\in U\),
\[
 e_+^\perp=\Lambda_{24}\oplus\mathbb Ze_+,
\qquad
 e_+^\perp/\mathbb Ze_+\cong\Lambda_{24}.
\]
La disminución de rango en dos procede de imponer ortogonalidad y cocientar
después por la recta nula. No es la reducción \(11\to10\) repetida con otros
números.

\section{Involución entre dato residual y cociente euclídeo}

La descomposición \(n=9Q+r\) sugiere dos coordenadas de naturaleza diferente:
el representante residual \(r\) y el cociente \(Q\). Para formular su
intercambio en un dominio estable, consideramos
\[
 \mathcal D=\mathbb Z^2\times\mathbb R_{>0},
\qquad
 E_R(m,w)=\frac{m^2}{R^2}+w^2R^2,
\]
y definimos
\[
 \mathsf T(m,w,R)=(w,m,R^{-1}).
\]
\begin{proposition}[Dualidad residual--euclídea]
\(\mathsf T\) es una involución y
\[
 E_R(m,w)=E_{R^{-1}}(w,m).
\]
Sobre \(\ell^2(\mathbb Z^2)\), el unitario
\(U_T|m,w\rangle=|w,m\rangle\) satisface
\[
 U_TH(R)U_T^{-1}=H(R^{-1}),
\qquad
 H(R)|m,w\rangle=E_R(m,w)|m,w\rangle.
\]
\end{proposition}
\begin{proof}
La segunda aplicación intercambia de nuevo las coordenadas e invierte dos
veces \(R\). La identidad de energía se obtiene por sustitución directa; la
identidad operatorial se comprueba sobre la base ortonormal.
\end{proof}

La proposición es una dualidad reticular exacta. La lectura física como
momento, enrollamiento y radio de compactificación requiere un mapa adicional
hacia el espacio de estados de una teoría de cuerdas; no se deduce de la
involución algebraica \cite{Giveon1994}.

\section{Interfaz tipada con la reducción undecadimensional}

La pantalla \(V_{11}=\ell_K\oplus V_{10}\) y la involución
\((m,w,R)\mapsto(w,m,R^{-1})\) proporcionan los dos datos algebraicos que
motivan la comparación con el contexto de teoría M: una reducción marcada
de once a diez coordenadas y una equivalencia de radios recíprocos. Para que
esa comparación sea una realización física debe especificarse un morfismo
\[
 \mathcal R_M:
 (\mathcal X_{\rm TPK},V_{11},\ell_K,\mathsf T)
 \longrightarrow
 (\mathcal H_M,\mathcal F_{10},R_{\rm phys},\alpha',
 \mathcal Q_{\rm susy}),
\]
donde el codominio incluye un espacio de estados, campos
decadimensionales, radio físico, tensión de cuerda y cargas supersimétricas
\cite{Witten1995}. La descomposición algebraica fija el dominio de
\(\mathcal R_M\); no introduce por sí sola osciladores, branas,
supersimetría ni las ecuaciones de supergravedad.

Esta separación conserva la profundidad de la correspondencia sin alterar
los tipos. La reducción \(11\to10\), la involución de radios y el cociente
\(26\to24\) son teoremas de los objetos definidos. La afirmación de que una
teoría dinámica concreta es su imagen exige exhibir \(\mathcal R_M\) y probar
que entrelaza acciones, restricciones y observables.

\section{Estado dimensional y transporte simultáneo}

El estado sobre el que actúa la Mecánica Dimensional se escribe
\[
 \mathfrak x=(x,\Sigma_n(x),\beta_n(x),[\gamma]),
\]
donde \(x\in\mathcal X_{\rm TPK}\), \(\Sigma_n(x)\) es su pantalla,
\(\beta_n(x)\) es el dato de frontera y \([\gamma]\) es la clase de trayectoria
compatible con ese dato. Un transporte no actúa sólo sobre \(x\): consta de
\[
 (T_\gamma,F_\gamma):
 (x,\Sigma_n,\beta_n)\longrightarrow
 (T_\gamma x,\Sigma_m,\beta_m),
\]
donde \(F_\gamma:M_n(x)\to M_m(T_\gamma x)\) envía relaciones en relaciones
y satisface el diagrama de compatibilidad con \(\beta_n,\beta_m\). El salto
dimensional es el cambio de rango del morfismo inducido
\[
 \overline F_\gamma:
 M_n(x)/\mathcal R_n(x)\longrightarrow
 M_m(T_\gamma x)/\mathcal R_m(T_\gamma x).
\]
Esta formulación evita atribuir una dimensión al punto aislado: la dimensión
pertenece al sistema de canales que acompaña al estado y a su transporte.

\section{Descenso a una dinámica sobre estados puntuales}

Sea \(q_{\rm pt}:\mathcal X_{\rm TPK}\to Z\) una proyección que elimina la
genealogía del estado, y definamos
\[
 Y_{\rm pt}:=\operatorname{im}(q_{\rm pt}).
\]
Sea, además, \(U:\mathcal X_{\rm TPK}\to\mathcal X_{\rm TPK}\) una
evolución completa.

\begin{theorem}[Criterio de descenso por fibras]
Existe una única aplicación
\(\overline U:Y_{\rm pt}\to Y_{\rm pt}\) tal que
\(q_{\rm pt}U=\overline Uq_{\rm pt}\) si y sólo si
\[
 q_{\rm pt}(x)=q_{\rm pt}(y)
 \Longrightarrow q_{\rm pt}(Ux)=q_{\rm pt}(Uy).
\]
\end{theorem}
\begin{proof}
La necesidad resulta aplicando \(\overline U\) a dos representantes de una
misma fibra. Para la suficiencia, se define
\[
 \overline U(q_{\rm pt}(x))=q_{\rm pt}(Ux).
\]
La condición garantiza que la definición no depende del representante. Como
el codominio se ha fijado igual a la imagen de \(q_{\rm pt}\), la
sobreyectividad requerida para la unicidad es automática.
\end{proof}

Cuando la condición falla, una misma clase puntual posee imágenes distintas;
la evolución descendida es entonces multivaluada. Conservar el dato de
memoria, restringir el dominio o trabajar con una correspondencia son tres
modificaciones matemáticas diferentes.

\section{Funcional aditivo sobre trayectorias orientadas}

Para evitar un término oculto de composición de trayectorias, el estado de una arista se
amplía con la dirección y el tipo aritmético de la arista precedente. Si
\(\widetilde e_j=(e_j,d_{j-1},t_{j-1})\), definimos los pesos locales
\[
 w_{\rm giro}(\widetilde e_j)=\mathbf1_{d_j\ne d_{j-1}},
 \qquad
 w_{\rm tipo}(\widetilde e_j)=\mathbf1_{t_j\ne t_{j-1}}.
\]
Los datos \((d_0,t_0)\) forman parte de la condición inicial de frontera.
Para una trayectoria aumentada
\(\widetilde\gamma=(\widetilde e_1,\ldots,\widetilde e_R)\), sean
\[
 N_{\rm giro}(\widetilde\gamma)=\sum_{j=1}^R
 w_{\rm giro}(\widetilde e_j),
 \qquad
 N_{\rm tipo}(\widetilde\gamma)=\sum_{j=1}^R
 w_{\rm tipo}(\widetilde e_j).
\]
Definimos
\[
 \mathbf I(\widetilde\gamma)=
 (N_{\rm giro}(\widetilde\gamma),N_{\rm tipo}(\widetilde\gamma)).
\]
La composición sólo se permite cuando el dato precedente del segundo camino
coincide con el dato terminal del primero. Con esta convención, la partición
de la suma en los dos intervalos de aristas demuestra
\[
 \mathbf I(\widetilde\gamma\widetilde\eta)
 =\mathbf I(\widetilde\gamma)+\mathbf I(\widetilde\eta),
\]
de modo que \(\mathbf I\) es un cociclo con valores en \(\NN^2\). Para
\(\lambda>0\), la combinación
\[
 I_\lambda(\widetilde\gamma)
 =N_{\rm giro}(\widetilde\gamma)
 +\lambda N_{\rm tipo}(\widetilde\gamma)
\]
es un funcional escalar sobre el espacio de trayectorias.

En la formulación lagrangiana ordinaria, el funcional de acción se define
sobre curvas de un espacio de configuración previamente elegido. Aquí el
orden es anterior: \(I_\lambda\) cuantifica cambios de dirección y de ley
aritmética en la propia categoría de trayectorias aumentadas. Un mapa de
realización puede transportar posteriormente este cociclo hacia una acción
continua. La aditividad demostrada arriba es la propiedad que permite esa
comparación sin identificar de antemano los dos funcionales.

\begin{principle}[Interacción mínima]
Entre trayectorias admisibles con extremos y datos de frontera fijados, la
evolución seleccionada minimiza \(I_\lambda\).
\end{principle}

\begin{proposition}[Existencia finita de minimizadores]
Si \(\Gamma(a,b)\) es finito y no vacío, \(I_\lambda\) alcanza un mínimo en
\(\Gamma(a,b)\).
\end{proposition}
\begin{proof}
La imagen de un conjunto finito no vacío por una función real posee un
elemento mínimo.
\end{proof}

El carácter de masa del \cref{chap:masas} constituye otro funcional sobre
trayectorias: transforma un cociclo entero en una razón positiva. El
funcional y el mecanismo que selecciona la trayectoria son objetos distintos;
la composición de ambos se especificará antes de introducir una unidad de
masa.
